% hw1-prop-logic.tex

\section{作业 1：命题逻辑}

\begin{problem}[命题逻辑的自然推理系统]
  请使用自然推理规则证明以下结论:

  {提示: 你可以使用 \LaTeX{} package `logicproof' 书写漂亮的证明。}
  \begin{enumerate}[(1)]
    \setcounter{enumi}{4}
    \item
      \[
        \vdash (\alpha \to (\beta \to \gamma)) \to ((\alpha \to \beta) \to (\alpha \to \gamma))
      \]
  \end{enumerate}
\end{problem}

\begin{solution}
  \begin{enumerate}[(1)]
    \setcounter{enumi}{4}
    \item
      \begin{intention}
        这是 A2 公理式; 依次假设
        $\alpha \to (\beta \to \gamma)$,
        $\alpha \to \beta$ 与 $\alpha$, 然后连做三次 MP 型推理。
      \end{intention}

      \begin{proof}
        \begin{logicproof}{3}
          \begin{subproof}
            \alpha \to (\beta \to \gamma) & assumption \\
            \begin{subproof}
              \alpha \to \beta & assumption \\
              \begin{subproof}
                \alpha & assumption \\
                \beta \to \gamma & $\to\elim$ 1, 3 \\
                \beta & $\to\elim$ 2, 3 \\
                \gamma & $\to\elim$ 4, 5
              \end{subproof}
              \alpha \to \gamma & $\to\intro$ 3--6
            \end{subproof}
            (\alpha \to \beta) \to (\alpha \to \gamma) & $\to\intro$ 2--7
          \end{subproof}
          (\alpha \to (\beta \to \gamma)) \to ((\alpha \to \beta) \to (\alpha \to \gamma)) & $\to\intro$ 1--8
        \end{logicproof}
      \end{proof}

      \begin{remark}
        \textbf{重点.} 在最内层假设 $\alpha$ 后, 先推出 $\beta \to \gamma$, 再推出 $\beta$, 最后得到 $\gamma$。
        \textbf{易错点.} 不要把第 4 行和第 5 行的先后关系写反: 必须先得到 $\beta \to \gamma$, 才能与 $\beta$ 合用。
      \end{remark}

  \end{enumerate}

\end{solution}

%%%%%%%%%%%%%%%

\begin{problem}[命题逻辑的自然推理系统]
  请使用自然推理规则证明以下结论:

  {提示: 你可以使用 \LaTeX{} package `logicproof' 书写漂亮的证明。}
  \begin{enumerate}[(1)]
    \setcounter{enumi}{5}
    \item
      \[
        \vdash (\lnot \beta \to \lnot \alpha) \to ((\lnot \beta \to \alpha) \to \beta)
      \]
  \end{enumerate}
\end{problem}

\begin{solution}
  \begin{enumerate}[(1)]
    \setcounter{enumi}{5}
    \item
      \begin{intention}
        先假设两条前提蕴含, 再以 $\lnot\beta$ 为临时假设导出矛盾, 从而得到 $\beta$。
      \end{intention}

      \begin{proof}
        \begin{logicproof}{4}
          \begin{subproof}
            \lnot\beta \to \lnot\alpha & assumption \\
            \begin{subproof}
              \lnot\beta \to \alpha & assumption \\
              \begin{subproof}
                \lnot\beta & assumption \\
                \lnot\alpha & $\to\elim$ 1, 3 \\
                \alpha & $\to\elim$ 2, 3 \\
                \bot & $\lnot\elim$ 4, 5
              \end{subproof}
              \lnot\lnot\beta & $\lnot\intro$ 3--6 \\
              \beta & $\lnot\lnot\elim$ 7
            \end{subproof}
            (\lnot\beta \to \alpha) \to \beta & $\to\intro$ 2--8
          \end{subproof}
          (\lnot\beta \to \lnot\alpha) \to ((\lnot\beta \to \alpha) \to \beta) & $\to\intro$ 1--9
        \end{logicproof}
      \end{proof}

      \begin{remark}
        \textbf{重点.} 在假设 $\lnot\beta$ 下, 两个前提分别给出 $\lnot\alpha$ 与 $\alpha$, 从而得到矛盾。
        \textbf{难点.} 由矛盾推出的不是直接的 $\beta$, 而是先得到 $\lnot\lnot\beta$, 再消去双重否定。
        \textbf{易错点.} 第 6 行应由 $\alpha$ 与 $\lnot\alpha$ 得到 $\bot$, 不能直接写成“故 $\beta$”。
      \end{remark}
  \end{enumerate}

\end{solution}
